\subsection{Carrier compressions on the minimal branch}

\begin{corollary}[Two-point spectral compression of $\gamma$ and $\phibase$]
Let
\[
y_+:=\frac12,
\qquad
y_-:=-\frac13
\]
be the two distinct carrier eigenvalues singled out by the factorized carrier polynomial. Then
\[
\gamma=y_+-y_-,
\qquad
\phibase=\frac{y_++y_-}{\pi}.
\]
Here $y_++y_-$ is the unweighted Vieta sum of the two distinct carrier eigenvalues, not the trace on $E$.
\end{corollary}

\begin{proof}
By \Cref{eq:carrier-polynomial}, equivalently by the factorization
\[
6Y^2-Y-\mathbf 1=(2Y-\mathbf 1)(3Y+\mathbf 1),
\]
the two distinct carrier eigenvalues are exactly $y_+=1/2$ and $y_-=-1/3$. Their difference is
\[
y_+-y_-=\frac12+\frac13=\frac56=\gamma,
\]
while their sum is
\[
y_++y_-=\frac12-\frac13=\frac16.
\]
Since $\phibase=(1-\gamma)/\pi=1/(6\pi)$, the claim follows.
\end{proof}

\begin{proposition}[Dual-root generation of the hypercharge packet and admissible cusp set]
\label{prop:dual-root-generation}
Let
\[
y_+=\frac12,
\qquad
y_-=-\frac13.
\]
Then all one-family hypercharges are generated by the same two carrier roots:
\[
\begin{aligned}
Y(Q_L)&=y_++y_-,
&
Y(u^c)&=2y_-,
&
Y(d^c)&=-y_-,
\\
Y(L_L)&=-y_+,
&
Y(e^c)&=2y_+,
&
Y(\nu^c)&=0.
\end{aligned}
\]
Moreover the positive transport cusps satisfy
\[
\left\{1,\frac23,\frac13\right\}=\{2y_+,-2y_-,2(y_++y_-)\}.
\]
Thus the carrier packet and the transport grammar are generated by the same two-point root data.
\end{proposition}

\begin{proof}
On the minimal carrier one has
\[
Y=\operatorname{diag}(y_-,y_-,y_-,y_+,y_+)
\]
with respect to the split $E_3\oplus E_2$. The one-family packet $S^+$ carries the sectors
$Q_L,u^c,d^c,L_L,e^c,\nu^c$, and their hypercharges are the sums of the occupied carrier roots.
This gives the displayed formulas. Substituting $y_+=1/2$ and $y_-=-1/3$ recovers the canonical
values
\[
\frac16,\quad
-\frac23,\quad
\frac13,\quad
-\frac12,\quad
1,\quad
0.
\]
The admissible transport cusps are fixed in \Cref{thm:exact-transport-closure} as
$\{1,\frac23,\frac13\}$, which is exactly the displayed root-generated set.
\end{proof}

\begin{corollary}[Family hypercharge variance and electroweak block trace]
Equip one carrier family with the uniform average
\[
\mathbb E_{S^+}(A):=\frac{1}{16}\tr_{S^+}A.
\]
Then
\[
\operatorname{Var}_{S^+}(Y)
:=
\mathbb E_{S^+}(Y^2)-\mathbb E_{S^+}(Y)^2
=
\frac{1}{16}\tr_{S^+}Y^2
=
\frac{5}{24}.
\]
Consequently,
\[
I_1^{\mathrm{EW}}=\operatorname{Var}_{S^+}(Y)\,\bone.
\]
\end{corollary}

\begin{proof}
The one-family trace identities give $\tr_{S^+}Y=0$ and $\tr_{S^+}Y^2=10/3$. Hence the centered second moment is
\[
\frac{1}{16}\tr_{S^+}Y^2=\frac{10/3}{16}=\frac{5}{24}.
\]
Using $I_1^{\mathrm{EW}}=\frac{5}{24}\,\bone$ from the electroweak block relation yields the second identity.
\end{proof}

\begin{corollary}[Factorial quadratic trace and two-defect factorization]
\label{cor:factorial-trace-delta2}
Let
\[
X:=6Y.
\]
Then the one-family quadratic trace is
\[
\tr_{S^+}(X^2)
=
36\,\tr_{S^+}(Y^2)
=
36\cdot\frac{10}{3}
=
120
=
5!.
\]
Equivalently, using the six hypercharge sectors with multiplicities
\[
(1,3,6,1,2,3)
\]
at
\[
X\in\{0,-4,1,6,-3,2\},
\]
the same trace is
\[
0^2+3\cdot 4^2+6\cdot 1^2+1\cdot 6^2+2\cdot 3^2+3\cdot 2^2=120.
\]
On the retained branch this gives the exact factorization
\[
\delta_2
=
2880\,\cthree^8
=
4!\,\tr_{S^+}(X^2)\,\cthree^8.
\]
\end{corollary}

\begin{proof}
By \Cref{prop:split-weight-enumerator}, one has
\[
\tr_{S^+}(Y^2)=\frac{10}{3}.
\]
Since $X=6Y$, the first identity follows immediately. The sector-by-sector sum is the same quadratic trace written in the basis of admissible occupancy sectors from \Cref{prop:family-hypercharge-compression}. Finally, the hard kernel normalization gives
\[
\delta_2=2880\,\cthree^8=24\cdot 120\,\cthree^8
=
4!\,\tr_{S^+}(X^2)\,\cthree^8.
\]
\end{proof}

\begin{remark}[Quadratic trace versus variance]
The value $120=5!$ is the unnormalized quadratic trace $\tr_{S^+}(X^2)$, not the normalized variance. The corresponding uniform variance is
\[
\operatorname{Var}_{S^+}(X)=36\,\operatorname{Var}_{S^+}(Y)=36\cdot\frac{5}{24}=\frac{15}{2}.
\]
\end{remark}

\begin{remark}[Generic abelian-index identity]
Define the generic occupancy
\[
\Omega_{\mathrm{occ}}:=16\,N_{\mathrm{fam}}.
\]
Then the abelian-index formula gives
\[
10\,\bone
=
\frac{40}{3}N_{\mathrm{fam}}+N_\Phi
=
\Omega_{\mathrm{occ}}\,\gamma+N_\Phi.
\]
At this stage no canonical specialization is inserted yet. The closed $41$ chain will be derived
only after the family and Higgs theorems fix $\Omega_{\mathrm{occ}}=48$ and $N_\Phi=1$.
\end{remark}

\begin{corollary}[Family integrality repair]
\label{cor:family-integrality-repair}
One carrier family contributes the weighted transport budget
\[
\dim S^+\cdot\gamma
=
16\cdot\frac56
=
\frac{40}{3}.
\]
Hence any integral global budget assembled from whole carrier families must satisfy
\[
N_{\mathrm{fam}}\cdot \frac{40}{3}\in\mathbb Z.
\]
The minimal positive repair is therefore
\[
N_{\mathrm{fam,min}}=3.
\]
For this minimal repair one gets
\[
\Omega_{\mathrm{occ}}=N_{\mathrm{fam,min}}\dim S^+=48,
\qquad
\Omega_{\mathrm{occ}}\gamma=40.
\]
The later family theorem realizes exactly this minimal arithmetic.
\end{corollary}

\begin{proof}
The first identity uses $\dim S^+=16$ and $\gamma=5/6$. Since $\gcd(40,3)=1$, integrality forces
$3\mid N_{\mathrm{fam}}$, so the smallest positive choice is $3$. Multiplying by $16$ and then by
$\gamma$ gives the displayed values.
\end{proof}

\subsection{Generic abelian index identity}

The abelian coefficient is first carried only in the generic form
\begin{equation}
\label{eq:41-chain-main}
10\,\bone
=
\Omega_{\mathrm{occ}}\gamma+N_\Phi.
\end{equation}
The arithmetic specialization to $41$ is deferred until the family multiplicity and the Higgs
sector have both been theorematically closed.

